problem:  
Let \((M,g,e^{-V}\mathrm{vol}_g)\) be a smooth, complete weighted Riemannian manifold satisfying the curvature–dimension condition \(CD(K,N)\) with \(K>0\) and finite \(N>1\).  
Let \(P_t = e^{t\Delta_V}\) denote the heat semigroup associated to \(\Delta_V = \Delta - \langle\nabla V,\nabla \cdot\rangle\).  
For \(x,y\in M\), define  
\[
\Phi_t(x,y) := W_2^2(P_t\delta_x,P_t\delta_y),
\]
the squared \(L^2\)-Wasserstein distance between heat-regularized Dirac measures. Under \(CD(K,N)\),
\[
\Phi_t(x,y)\le e^{-2Kt}d^2(x,y)
\]
for small \(t>0\), with equality in Gaussian space.

Suppose that, for small \(t>0\), one has the **uniform first-term expansion**
\[
\frac{\Phi_t(x,y)}{d^2(x,y)} = e^{-2Kt}\bigl(1-\alpha\,t + o(t)\bigr),\quad t\to0,
\]
where \(\alpha\ge0\) is a constant independent of \(x,y\).  

**Problem:**  
Determine whether the requirement that the first-order deviation coefficient \(\alpha\) be spatially constant forces the Bakry–Émery tensor to be constant, i.e.
\[
\operatorname{Ric}_V \equiv K g,
\]
and hence \((M,g,V)\) to be Gaussian up to affine equivalence.  
Equivalently, can any nontrivial \(CD(K,N)\) weighted manifold exhibit uniform small-time diffusion-distance decay independent of base points?  
More generally, describe how the small-time expansion of \(\Phi_t(x,y)\) determines local geometric data (e.g., curvature, Hessian of \(V\)), and whether uniformity of these coefficients implies rigidity.

Why is it a "good" problem:  
This problem transforms global equality questions in the heat-flow contraction principle into a *local asymptotic rigidity problem*, probing how fine-scale diffusion–transport behavior encodes curvature. By focusing on the coefficient of the first deviation from exponential decay, it directly links the stochastic geometry of Wasserstein distances to local Ricci–Bakry–Émery invariants.  
If spatial constancy of this coefficient necessarily enforces constant weighted Ricci curvature, it would establish a new route to identify Gaussian structures from infinitesimal transport data—bridging heat kernel asymptotics, differential geometry, and optimal transport.  
The formulation is simple but deep: it demands understanding of how small-time transport asymptotics reflect curvature-dimension conditions and whether any non-Gaussian geometry can present globally uniform diffusion symmetry. Resolving it would sharpen our grasp of rigidity phenomena at the interface of analysis and geometry in the \(CD(K,N)\) framework.